\documentclass[11pt,a4paper]{article}

\usepackage[french]{babel}
\usepackage[T1]{fontenc}
\usepackage[utf8]{inputenc}
\usepackage{lmodern}
\usepackage{geometry}
\usepackage{xcolor}
\usepackage{longtable}
\usepackage{booktabs}
\usepackage{hyperref}
\usepackage{listings}
\usepackage{parskip}
\usepackage{array}

\geometry{margin=2.5cm}

\hypersetup{
	colorlinks=true,
	linkcolor=blue!70!black,
	urlcolor=blue!70!black,
}

\definecolor{codebg}{gray}{0.96}
\definecolor{codefg}{rgb}{0.1,0.1,0.4}
\definecolor{commentfg}{rgb}{0.4,0.4,0.4}

\lstset{
	language=[LaTeX]TeX,
	basicstyle=\ttfamily\small,
	keywordstyle=\color{codefg}\bfseries,
	commentstyle=\color{commentfg}\itshape,
	backgroundcolor=\color{codebg},
	frame=single,
	framerule=0.4pt,
	rulecolor=\color{gray!50},
	xleftmargin=1em,
	xrightmargin=1em,
	aboveskip=0.8em,
	belowskip=0.8em,
	columns=fullflexible,
	keepspaces=true,
	showstringspaces=false,
	breaklines=true,
	morekeywords={
		teiP,teiHi,teiHead,teiNote,teiRef,teiQ,teiTitle,teiName,
		teiPersName,teiPlaceName,teiDate,teiNum,teiOrgName,teiTerm,
		teiSaid,teiGraphic,teiLb,teiPb,teiPtr,teiLabel,teiItem,
		teiAuthor,teiBiblScope,teiEditor,teiForeign,teiIdno,
		teiPublisher,
		lateiNoIndent,lateiIndent,
		lateiVSpace,lateiSpaceBefore,lateiSpaceAfter,
		lateiPageBreak,lateiClearPage,
		lateiPageBreakBefore,lateiPageBreakAfter,
		lateiClearPageBefore,lateiClearPageAfter,
		lateiKeepWithNext,lateiKeepTogether,
		lateiNoPageBreakBefore,lateiNoPageBreakAfter,
	},
}

% Commande pour afficher un terme technique
\newcommand{\cmd}[1]{\texttt{\textbackslash #1}}
\newcommand{\attrib}[1]{\texttt{#1}}
\newcolumntype{L}[1]{>{\raggedright\arraybackslash}p{#1}}
% ============================================================
\title{%
	\textbf{LaTEI : guide des corrections éditoriales de mise en page}\\[0.4em]
}
\author{PURH — Usage interne — Impression}
\date{}
% ============================================================

\begin{document}
	
	\maketitle
	\tableofcontents
	\newpage
	
	% ============================================================
	\section{Principe général}
	% ============================================================
	
	Le logiciel Impression est un générateur de livres savants qui produit du html statique ou du PDF à partir de fichiers XML Métopes. Il est téléchargeable sous licence libre ici: https://github.com/Gheeraert/impression
	
	Impression produit un LaTeX modifiable afin d'être réengistré au format d'origine XML Métopes - Common Publishing. Pour y parvenir, il mobilise une syntaxe particulière à base de commandes LaTeX, le LaTEI.
	
	Un fichier \texttt{.latei.tex} contient deux types de commandes qu'il est
	important de distinguer.
	
	\subsection{Les trois couches du LaTEI}
	
	\begin{description}
		\item[\cmd{tei...}] Commandes éditoriales réversibles. Elles encodent le
		contenu du livre~: textes, titres, notes, figures, liens, tableaux.
		Lors de la restauration XML/Métopes, \emph{ces commandes sont
			retranscrites telles quelles} vers le XML d'origine.
		
		\item[\cmd{latei...}] Corrections locales de mise en page PDF. Elles
		servent uniquement à améliorer l'apparence du PDF~: supprimer un
		alinéa, ajouter un blanc vertical, forcer un saut de page, éviter
		une coupure indésirable. Lors de la restauration XML, \emph{ces
			commandes disparaissent complètement}~; seul leur contenu éditorial
		est conservé.
		
		\item[LaTeX brut] Commandes LaTeX ordinaires comme \cmd{noindent},
		\cmd{vspace}, \cmd{newpage}, \cmd{clearpage}. Elles sont
		\textbf{interdites dans la zone éditoriale}. Si vous en écrivez,
		la restauration XML sera bloquée avec une erreur.
	\end{description}
	
	\subsection{La règle à retenir}
	
	\begin{itemize}
		\item Les commandes \cmd{tei...} portent le \emph{contenu}~: texte,
		structure, métadonnées.
		\item Les commandes \cmd{latei...} portent uniquement la \emph{mise en
			page PDF}~; elles n'existent pas dans le XML.
		\item Toute commande non documentée dans ce guide est interdite dans la
		zone éditoriale.
	\end{itemize}
	
	% ============================================================
	\section{Ce qu'on peut modifier}
	% ============================================================
	
	\subsection{Modifications autorisées}
	
	Les éditrices et éditeurs peuvent librement~:
	
	\begin{itemize}
		\item corriger une coquille ou une faute de frappe dans le texte~;
		\item ajouter ou supprimer un mot~;
		\item modifier la ponctuation, une espace, une apostrophe, une
		majuscule~;
		\item corriger une espace insécable ou un tiret~;
		\item ajouter une commande \cmd{latei...} documentée pour corriger la
		mise en page du PDF.
	\end{itemize}
	
	\textbf{Exemple autorisé} — corriger un mot dans le texte~:
	
	\begin{lstlisting}
		% Avant correction :
		\teiP[xmlid={p-111}]{C'est effectivement ce qui s'était pratiqué.}
		
		% Apres correction :
		\teiP[xmlid={p-111}]{C'est bien ce qui s'était pratiqué.}
	\end{lstlisting}
	
	\subsection{Modifications interdites}
	
	Il ne faut pas~:
	
	\begin{itemize}
		\item renommer ou abréger une commande \cmd{tei...}~;
		\item supprimer ou déplacer une accolade~;
		\item déplacer une commande \cmd{tei...} dans la structure~;
		\item supprimer ou modifier un attribut comme \attrib{xmlid=\{...\}}~;
		\item écrire du LaTeX brut non documenté~;
		\item transformer une correction de contenu en correction de mise en page.
	\end{itemize}
	
	\textbf{Exemples interdits}~:
	
	\begin{lstlisting}
		% Interdit - commande tei renommee :
		\p[xmlid={p-111}]{Texte.}
		
		% Interdit - LaTeX brut :
		\noindent\teiP[xmlid={p-111}]{Texte.}
		
		% Interdit - dimension libre :
		\lateiVSpace{12pt}
	\end{lstlisting}
	
	% ============================================================
	\section{Liste officielle des commandes de mise en page}
	% ============================================================
	
	Toutes les commandes ci-dessous sont transparentes lors de la restauration
	XML~: elles modifient le rendu du PDF, mais le XML restauré ne les connaît
	pas.
	
	\subsection{Alinéas}
	
	\begin{lstlisting}
		\lateiNoIndent{...}   % supprime le retrait de premiere ligne
		\lateiIndent{...}     % force le retrait de premiere ligne
	\end{lstlisting}
	
	\subsection{Espacements verticaux}
	
	Les valeurs autorisées pour la taille sont~: \texttt{small} (environ un
	demi-interligne), \texttt{medium} (environ un interligne), \texttt{large}
	(environ deux interlignes). Toute autre valeur est refusée.
	
	\begin{lstlisting}
		% Blanc vertical autonome (entre deux elements) :
		\lateiVSpace{small}
		\lateiVSpace{medium}
		\lateiVSpace{large}
		
		% Blanc avant le contenu :
		\lateiSpaceBefore{small}{...}
		\lateiSpaceBefore{medium}{...}
		\lateiSpaceBefore{large}{...}
		
		% Blanc apres le contenu :
		\lateiSpaceAfter{small}{...}
		\lateiSpaceAfter{medium}{...}
		\lateiSpaceAfter{large}{...}
	\end{lstlisting}
	
	\subsection{Sauts de page}
	
	\begin{lstlisting}
		% Commandes autonomes (entre deux elements) :
		\lateiPageBreak        % saut de page simple
		\lateiClearPage        % saut de page + vidage des figures en attente
		
		% Commandes enveloppantes :
		\lateiPageBreakBefore{...}    % saut avant le contenu
		\lateiPageBreakAfter{...}     % saut apres le contenu
		\lateiClearPageBefore{...}    % clearpage avant le contenu
		\lateiClearPageAfter{...}     % clearpage apres le contenu
	\end{lstlisting}
	
	\subsection{Contrôle des coupures de page}
	
	\begin{lstlisting}
		\lateiKeepWithNext{...}       % interdit la coupure apres ce contenu
		\lateiKeepTogether{...}       % maintient tout le contenu sur une page
		\lateiNoPageBreakBefore{...}  % interdit la coupure avant ce contenu
		\lateiNoPageBreakAfter{...}   % interdit la coupure apres ce contenu
	\end{lstlisting}
	
	% ============================================================
	\section{Exemples pratiques}
	% ============================================================
	
	\subsection{Ajouter un espace entre deux paragraphes}
	
	\begin{lstlisting}
		\teiP[xmlid={p-111}]{Premier paragraphe.}
		
		\lateiVSpace{small}
		
		\teiP[xmlid={p-112}]{Second paragraphe avec un peu d'espace au-dessus.}
	\end{lstlisting}
	
	\subsection{Supprimer un alinéa}
	
	\begin{lstlisting}
		\lateiNoIndent{
			\teiP[xmlid={p-112}]{Ce paragraphe n'a pas d'alinea.}
		}
	\end{lstlisting}
	
	\subsection{Ajouter un blanc avant un paragraphe}
	
	\begin{lstlisting}
		\lateiSpaceBefore{medium}{
			\teiP[xmlid={p-112}]{Ce paragraphe a un blanc avant lui.}
		}
	\end{lstlisting}
	
	\subsection{Forcer un saut de page}
	
	\begin{lstlisting}
		\teiP[xmlid={p-200}]{Fin d'une section.}
		
		\lateiPageBreak
		
		\teiP[xmlid={p-201}]{Debut de la page suivante.}
	\end{lstlisting}
	
	\subsection{Garder un titre avec le paragraphe suivant}
	
	\begin{lstlisting}
		\lateiKeepTogether{
			\teiHead{Un titre a ne pas isoler}
			\teiP[xmlid={p-300}]{Premier paragraphe sous le titre.}
		}
	\end{lstlisting}
	
	\subsection{Combiner plusieurs corrections}
	
	Il est possible d'imbriquer les commandes \cmd{latei...}. Par exemple,
	ajouter un blanc avant un paragraphe sans alinéa~:
	
	\begin{lstlisting}
		\lateiSpaceBefore{small}{
			\lateiNoIndent{
				\teiP[xmlid={p-112}]{Paragraphe sans alinea, avec un blanc avant.}
			}
		}
	\end{lstlisting}
	
	Le XML restauré ne contient que le paragraphe~:
	
	\begin{lstlisting}[language=XML]
		<p xml:id="p-112">Paragraphe sans alinea, avec un blanc avant.</p>
	\end{lstlisting}
	
	% ============================================================
	\section{Que devient le XML après correction ?}
	% ============================================================
	
	Les commandes \cmd{latei...} n'existent pas dans le XML. Lors de la
	restauration, elles sont ignorées~; seul leur contenu est restauré.
	
	\textbf{Exemple} — dans le fichier \texttt{.latei.tex}~:
	
	\begin{lstlisting}
		\lateiSpaceBefore{small}{
			\teiP[xmlid={p-112}]{Texte corrige.}
		}
	\end{lstlisting}
	
	\textbf{XML restauré}~:
	
	\begin{lstlisting}[language=XML]
		<p xml:id="p-112">Texte corrige.</p>
	\end{lstlisting}
	
	La commande \texttt{\textbackslash lateiSpaceBefore\{small\}\{...\}} a
	disparu. L'attribut \texttt{xml:id} et le texte sont intacts.
	
	\medskip
	\textbf{Ce qui ne change jamais dans le XML~:}
	\begin{itemize}
		\item le texte et sa structure~;
		\item les attributs (\attrib{xmlid}, \attrib{rend}, \attrib{type}, etc.)~;
		\item la hiérarchie des éléments (divisions, paragraphes, listes, etc.)~;
		\item les notes, les références, les images.
	\end{itemize}
	
	% ============================================================
	\section{Tableau récapitulatif}
	% ============================================================
	
	\small
	\setlength{\tabcolsep}{4pt}
	\begin{longtable}{@{}L{0.30\textwidth}L{0.22\textwidth}L{0.22\textwidth}L{0.18\textwidth}@{}}
		\toprule
		\textbf{Commande} & \textbf{Usage} & \textbf{Où la placer} & \textbf{Effet sur le XML} \\
		\midrule
		\endhead
		\midrule
		\multicolumn{4}{r}{\small Suite page suivante\ldots} \\
		\endfoot
		\bottomrule
		\endlastfoot
		
		\cmd{lateiNoIndent\{...\}}
		& Supprimer l'alinéa
		& Envelopper le paragraphe
		& Aucun (contenu conservé) \\[0.4em]
		
		\cmd{lateiIndent\{...\}}
		& Forcer l'alinéa
		& Envelopper le paragraphe
		& Aucun (contenu conservé) \\[0.4em]
		
		\cmd{lateiVSpace\{size\}}
		& Blanc vertical
		& Entre deux éléments
		& Ignorée \\[0.4em]
		
		\cmd{lateiSpaceBefore\{size\}\{...\}}
		& Blanc avant
		& Envelopper l'élément
		& Aucun (contenu conservé) \\[0.4em]
		
		\cmd{lateiSpaceAfter\{size\}\{...\}}
		& Blanc après
		& Envelopper l'élément
		& Aucun (contenu conservé) \\[0.4em]
		
		\cmd{lateiPageBreak}
		& Saut de page
		& Entre deux éléments
		& Ignorée \\[0.4em]
		
		\cmd{lateiClearPage}
		& Nouvelle page + flottants
		& Entre deux éléments
		& Ignorée \\[0.4em]
		
		\cmd{lateiPageBreakBefore\{...\}}
		& Saut avant
		& Envelopper l'élément
		& Aucun (contenu conservé) \\[0.4em]
		
		\cmd{lateiPageBreakAfter\{...\}}
		& Saut après
		& Envelopper l'élément
		& Aucun (contenu conservé) \\[0.4em]
		
		\cmd{lateiClearPageBefore\{...\}}
		& Clearpage avant
		& Envelopper l'élément
		& Aucun (contenu conservé) \\[0.4em]
		
		\cmd{lateiClearPageAfter\{...\}}
		& Clearpage après
		& Envelopper l'élément
		& Aucun (contenu conservé) \\[0.4em]
		
		\cmd{lateiKeepWithNext\{...\}}
		& Éviter coupure après
		& Envelopper l'élément
		& Aucun (contenu conservé) \\[0.4em]
		
		\cmd{lateiKeepTogether\{...\}}
		& Bloc sur une page
		& Envelopper le bloc
		& Aucun (contenu conservé) \\[0.4em]
		
		\cmd{lateiNoPageBreakBefore\{...\}}
		& Éviter coupure avant
		& Envelopper l'élément
		& Aucun (contenu conservé) \\[0.4em]
		
		\cmd{lateiNoPageBreakAfter\{...\}}
		& Éviter coupure après
		& Envelopper l'élément
		& Aucun (contenu conservé) \\
		
	\end{longtable}
	\normalsize
	
	\noindent\textbf{Tailles valides} pour \texttt{size}~: \texttt{small},
	\texttt{medium}, \texttt{large}. Toute autre valeur est refusée.
	
	% ============================================================
	\section{Document évolutif}
	% ============================================================
	
	Ce guide décrit les commandes actuellement disponibles. De nouvelles
	commandes pourront être ajoutées si des besoins récurrents apparaissent
	en correction d'épreuves.
	
	\textbf{Seules les commandes documentées ici doivent être utilisées.}
	Toute commande LaTeX non documentée dans ce guide sera refusée lors de
	la restauration XML.
	
	Si un besoin de correction récurrent n'est pas couvert par les commandes
	ci-dessus, il convient de le signaler pour qu'une commande dédiée soit
	ajoutée officiellement, plutôt que d'écrire du LaTeX brut.
	
\end{document}
